\documentclass[10pt,a4paper]{amsart}
\usepackage[margin=19mm]{geometry}
\usepackage{amsmath,amssymb,mathtools}
\usepackage[scr=rsfs,cal=euler]{mathalfa}
\usepackage{xfrac}
\usepackage[unicode,hidelinks,bookmarks=false]{hyperref}
\newcommand{\ie}{i.e.\ }
\newcommand{\cf}{cf.\ }
\newcommand{\resp}{respectively\ }
\newcommand{\Subsection}{\subsection}
\providecommand{\nobreakdash}{\nobreak\hbox{-}}
\setlength{\emergencystretch}{2em}
\multlinegap0pt
\usepackage{bm}
\usepackage{bbm}
\usepackage{dsfont}
\usepackage{xspace}
\usepackage{cancel}
\usepackage{mathtools}
\usepackage{amsthm}
\makeatletter
\@ifundefined{definition}{\newtheorem{definition}{Definition}}{}
\@ifundefined{lemma}{\newtheorem{lemma}[definition]{Lemma}}{}
\@ifundefined{proposition}{\newtheorem{proposition}[definition]{Proposition}}{}
\makeatother
\makeatletter
\newcommand{\W}{\mathsf{W}}
\def\Z{\mathbb{Z}}
\expandafter\ifx\csname customthm\endcsname\relax
\newenvironment{customthm}[1]
{\renewcommand\theinnercustomthm{#1}\innercustomthm}
\fi
\newcommand{\fine}{\mathsf{fine}}
\newcommand{\gr}{\mathrm{gr}}
\newcommand{\tors}{\mathrm{tors}}
\makeatother
\title[Genus one correspondence between tropical and algebraic curv]{Genus one correspondence between tropical and algebraic curves}
\author{Alessio Cela, Sae Koyama}
\date{Source: arXiv:2607.06426v2 (CC BY 4.0)}
\begin{document}
\maketitle

\noindent\emph{Source and licence.} Genus one correspondence between tropical and algebraic curves. Version arXiv:2607.06426v2 (CC BY 4.0), distributed under the Creative Commons Attribution 4.0 International licence, as recorded in the arXiv licence metadata for this exact version and shown on the abstract page https://arxiv.org/abs/2607.06426v2. Full rights notice: licenses/arxiv-2607.06426-CC-BY-4.0.txt.\par\smallskip

\noindent\textit{Editorial scope.} Counted unit: the Proposition whose source label is \texttt{prop: cycle induction} in the article (source0.tex lines 4830--4834, source label \texttt{prop: cycle induction}), delivered together with the article's own complete original proof of it (source0.tex lines 4836--4936, one \texttt{proof} environment of 101 lines). No mathematics is written, repaired or re-derived: statement and proof are the article's own text line for line. Required context retained uncounted: eqn: relations from tropical curve (lines 3484--3486); lem: constrained (lines 4675--4691). Every definition, display and label of those lines is retained unchanged. Omitted from this package and not redistributed: 1--3483, 3487--4674, 4692--4829, 4835--4835, 4937--5083. Nothing inside a retained excerpt is omitted or altered: every retained body is delivered exactly as the article's lines are written, so no omission receipt of any kind accompanies this package. Citations: the retained excerpts carry no citation command at all, so no bibliography entry of the article is transcribed and no reference is redistributed. Reference adaptation: none was needed and none was made; every reference inside the delivered unit resolves inside it, so no unresolved reference or citation is produced. Printed slips kept literal: no printed word, symbol, hypothesis or number of the counted statement or of its proof was altered. Layout and class: the article's own class is \texttt{article} with packages including geometry, titlesec, tikz, tikz-cd, url, hyperref, float, amsmath, stmaryrd, amssymb, enumitem, graphicx, mathdots, color; the wrapper is the lane's pinned \texttt{amsart} editorial wrapper at 10pt on A4, which restates the theorem-like environments the retained text needs and adds the article's own macro definitions that the retained text needs (\texttt{\textbackslash W}, \texttt{\textbackslash Z}, \texttt{\textbackslash fine}, \texttt{\textbackslash gr}, \texttt{\textbackslash tors}), with extra wrapper packages (bm, bbm, dsfont, xspace, cancel, mathtools). The delivered numbering is the unit's own. Assets: no retained span contains a figure environment, a graphics-inclusion command, a PDF-page inclusion, a table or a TikZ picture; no image file is copied into this package, the package declares no asset dependency and no image file is redistributed with it. Licence: version arXiv:2607.06426v2 of this article is distributed under the Creative Commons Attribution 4.0 International licence, as recorded in the arXiv licence metadata for this exact version and shown on the abstract page https://arxiv.org/abs/2607.06426v2. A full rights notice is delivered at licenses/arxiv-2607.06426-CC-BY-4.0.txt. Characters: every retained excerpt is pure ASCII, so the delivered engine, XeLaTeX with its default Latin Modern text font, reports no missing character.
        \begin{equation}\label{eqn: relations from tropical curve}
        a_{\vec{e}}(m) = \left((\ldots, \iota_{\vartheta_V}(m), \ldots, -\iota_{\vartheta_{V'}}(m), \ldots), (\ldots, u_{\vec{e}}(m), \ldots)\right),
        \end{equation}

\begin{lemma} \label{lem: constrained}
    Let $M = \mathbb{Z}^k \times M_{\mathrm{tors}}$ be an abelian group with torsion subgroup $M_{\mathrm{tors}}$. Fix a positive integer $d \in \mathbb{Z}_{>0}$, a homomorphism $\alpha: \mathbb{Z}^r \to M$, and integers $a_1, \ldots, a_{r-1}$. Consider the quotient group
\begin{equation*}
    N = (M \times \mathbb{Z}^{r-1} \times \mathbb{Z}) / R
\end{equation*}
where the relation $R$ is the subgroup generated by the image of the map
\begin{align*}
    \mathbb{Z}^r & \to M \times \mathbb{Z}^{r-1} \times \mathbb{Z} \\
    (m_1, \ldots, m_r) & \mapsto \left(\alpha(m), (m_1, \ldots, m_{r-1}), \sum_{i=1}^{r-1} a_i m_i + d m_r\right).
\end{align*}
Then,  $|M_{\mathrm{tors}} |$ divides $ |N_{\mathrm{tors}}|$. Let $g \in \Z_{\geq 1}$ be the result of this division. It follows that $g$ divides $d$ and that the induced homomorphism
\begin{equation*}
    \Phi: M/M_{\mathrm{tors}} \to N/N_{\mathrm{tors}}
\end{equation*}
is an isomorphism after tensoring with $\mathbb{R}$ and has lattice index equal to $d/g$.
    
\end{lemma}

\begin{proposition}\label{prop: cycle induction}
    For all $0<k \leq E_1$, the cardinality of the torsion subgroup of $Q^\fine_\W(\sigma^{(k-1)})^{\gr}$ divides that of the torsion subgroup of $Q^\fine_\W(\sigma^{(k)})^{\gr}$, and one has 
    $$  \frac{|(Q^\fine_\W(\sigma^{(k)})^{\gr})_{\tors}|}{|(Q^\fine_\W(\sigma^{(k-1)})^{\gr})_{\tors}|} = \mathrm{lat}(\sigma^{(k-1)})  \cdot \frac{ d_{V_k}}{\mathrm{lat}(\sigma^{(k)})}
    $$  
\end{proposition}

\begin{proof}
    Let $\vartheta:=\vartheta_{V_k}$ be the codimension $1$ subcone of $\Sigma(X)$ which contains $h^{(k)}(V_k)$. Let  $\vec{e}_1 \colon V_{i_1} \rightarrow V_k$ and $\vec{e}_2 \colon V_k \rightarrow V_{i_2}$ be the two directed edges adjacent to $V_k$. Let $\vartheta_1:= \vartheta_{V_{i_1}}, \vartheta_2=\vartheta_{V_{i_2}} $ be the cones of $\Sigma(X)$ containing $V_{i_1}, V_{i_2}$ (which may be of maximum dimension, or codimension $1$). Let $\iota_{\theta_i} :  M_X  \rightarrow M_{\vartheta_i}$ be dual to the associated inclusion of monoids $N_{\vartheta_i} \rightarrow N_X$. Let $\vec{e}_0 \colon V_{i_1} \rightarrow V_{i_2}$ be the directed edge in $\sigma^{(k-1)}$. We have 
    $$Q_{\W}^{\fine}(\sigma^{(k-1)})^{\gr} = \frac{(\prod_{V \neq V_{i_1}, V_{i_2}} M_{\vartheta_V} \times \prod_{e \neq e_0} \mathbb{Z}) \times M_{\vartheta_1} \times M_{\vartheta_2} \times  \mathbb{Z}_{\vec{e}_0}}{R' + ((0), -\iota_1(m), \iota_2(m), u_{\vec{e}_0}(m))},$$
    where we denoted by $R'$ the set of relations \eqref{eqn: relations from tropical curve} coming from edges in $\sigma^{(k-1)}$ different from $e_0$ and $m$ varies in $M_X$. Meanwhile  
    $$Q_{\W}^{\fine}(\sigma^{(k)})^{\gr} = \frac{(\prod_{V \neq V_{i_1}, V_{i_2}} M_{\vartheta_V} \times \prod_{e \neq e_0} \mathbb{Z}) \times M_{\vartheta_1} \times M_{\vartheta_2}  \times M_{\vartheta} \times \mathbb{Z}_{\vec{e}_1} \times \mathbb{Z}_{\vec{e}_2} }{R' + (a_{\vec{e}_1}(m)) + (a_{\vec{e}_2}(m))}.$$
    where for $m \in M_X$
    \begin{align*}
        &a_{\vec{e}_1}(m) = ((0), -\iota_1(m), 0 , \iota(m),  u_{\vec{e}_1}(m), 0)\\
        &a_{\vec{e}_2}(m) = ((0), 0, \iota_2(m) , -\iota(m), 0, u_{\vec{e}_2}(m))
    \end{align*}
    Note that $u_{\vec{e}_0} = u_{\vec{e}_1} = u_{\vec{e}_2}$. The two relations $a_{\vec{e}_1}(m)$ and $a_{\vec{e}_2}(m)$ can be replaced with 
    \begin{align*}
        &((0), -\iota_1(m), \iota_2(m),  0,  u_{\vec{e}_0}(m), u_{\vec{e}_0}(m))  \quad \text{and} \quad
        ((0), 0, \iota_2(m) , -\iota(m), 0, u_{\vec{e}_0}(m)). 
    \end{align*}
    Consider a change of basis $\mathbb{Z}_{\vec{e}_1} \times \mathbb{Z}_{\vec{e}_2} \rightarrow \mathbb{Z} \times \mathbb{Z}$ given by $(n,m) \mapsto (n, n-m)$. Under this change of basis, the two relations become 
    \begin{align*}
            &((0), -\iota_1(m), \iota_2(m),  0,  u_{\vec{e}_0}(m), 0)  \quad \text{and} \quad
            ((0), 0, \iota_2(m) , -\iota(m), 0, u_{\vec{e}_0}(m)). 
    \end{align*}
    The first relation is one of the defining relations of $Q_{\W}^{\fine}(\sigma^{(k-1)})^{\gr}$. Thus, we have an isomorphism
\[
Q_{\W}^{\fine}(\sigma^{(k)})^{\gr}
\simeq
\frac{
Q_{\W}^{\fine}(\sigma^{(k-1)})^{\gr}
\times M_{\vartheta}
\times \mathbb{Z}
}{
(\alpha(m),-\iota(m),u_{\vec e_0}(m))
}.
\]
Here
\[
\alpha:M_X\longrightarrow Q_\W^\fine(\sigma^{(k-1)})^{\gr}
\]
sends $m\in M_X$ to the image of
\[
((0),0,\iota_2(m),0)
\in
\Bigl(
\prod_{V\neq V_{i_1},V_{i_2}} M_V
\times
\prod_{e\neq e_0}\mathbb Z
\Bigr)
\times M_{\vartheta_1}
\times M_{\vartheta_2}
\times \mathbb Z_{\vec e_0}
\]
under the projection to $Q_\W^\fine(\sigma^{(k-1)})^{\gr}$.

Choose bases of $M_X$ and $M_{\vartheta}$ such that
\[
\iota_\vartheta:M_X\to M_\vartheta,
\qquad
(m_1,\ldots,m_r)\longmapsto (m_1,\ldots,m_{r-1}).
\]
Then, by Lemma~\ref{lem: constrained}, the natural map
\[
Q_\W^\fine(\sigma^{(k-1)})^\gr
\longrightarrow
Q_{\W}^{\fine}(\sigma^{(k)})^{\gr}
\simeq
\frac{
Q_{\W}^{\fine}(\sigma^{(k-1)})^{\gr}
\times M_{\vartheta}
\times \mathbb{Z}
}{
\bigl(
\alpha(m_1,\ldots,m_r),
-(m_1,\ldots,m_{r-1}),
\sum_{i=1}^{r-1} a_i m_i + d_{V_k}m_r
\bigr)
}
\]
becomes an isomorphism after tensoring with $\mathbb R$, and
\[
\frac{\bigl|(Q_\W^\fine(\sigma^{(k)})^{\gr})_{\tors}\bigr|}
     {\bigl|(Q_\W^\fine(\sigma^{(k-1)})^{\gr})_{\tors}\bigr|}
=
\frac{d_{V_k}}
{\left[
Q_\W^\fine(\sigma^{(k)})^{\gr}/(Q_\W^\fine(\sigma^{(k)})^{\gr})_{\tors}
:
\operatorname{im}\!\Bigl(
Q_\W^\fine(\sigma^{(k-1)})^{\gr}/(Q_\W^\fine(\sigma^{(k-1)})^{\gr})_{\tors}
\Bigr)
\right]}.
\]
The conclusion now follows from the identity
$$
\left[
Q_\W^\fine(\sigma^{(k)})^{\gr}/(Q_\W^\fine(\sigma^{(k)})^{\gr})_{\tors}
:
\operatorname{im}\!\Bigl(
Q_\W^\fine(\sigma^{(k-1)})^{\gr}/(Q_\W^\fine(\sigma^{(k-1)})^{\gr})_{\tors}
\Bigr)
\right]= \frac{ \mathrm{lat}(\sigma^{(k)})}{\mathrm{lat}(\sigma^{(k-1)})}
$$
    
\end{proof}

\end{document}
